\documentclass[11pt]{article}
\usepackage[margin=1in]{geometry}
\usepackage{amsmath,amssymb}
\usepackage{booktabs}
\usepackage{microtype}
\usepackage[hidelinks]{hyperref}
\usepackage{titlesec}
\titlespacing*{\section}{0pt}{1.1ex}{0.6ex}
\titlespacing*{\subsection}{0pt}{0.8ex}{0.4ex}
\setlength{\parskip}{0.35em}
\setlength{\parindent}{0pt}

\title{\vspace{-3em}\textbf{GIP: Generalized Intuition Pretraining in Padded\\
Joint-Embedding Predictive Architectures for Multi-Task Control}}
\author{Minghao Fu}
\date{\today}

\begin{document}
\maketitle
\vspace{-1.5em}

\begin{abstract}
\noindent
A world model can control in two ways that rarely meet: by \emph{searching} (model-predictive
control over imagined rollouts, accurate but slow and high-variance) or by \emph{reacting} (an
amortized policy, fast but myopic). We propose GIP, a latent world model that carries a fast
\emph{intuition head}, an amortized action predictor living in the model's own latent space, and
uses it to guide slow deliberative planning. The intuition proposes a short action sequence that
warm-starts the planner; the world model then rolls candidates to the goal and selects. Across
goal-reaching control this intuition-guided planning matches or beats both pure planning and the
reactive policy, and reaches the planner's ceiling in 4--5$\times$ fewer search iterations. We
pretrain the intuition jointly across seven heterogeneous tasks inside a single padded
joint-embedding predictive architecture (JEPA), using per-task action masks, a frozen language
task token, and a predicted proprioception token. We propose to scale this generalized intuition
pretraining toward a reusable multi-task control prior.
\end{abstract}

\section{Motivation}
A world model affords two modes of control. Model-predictive control, for instance the cross-entropy
method (CEM) over imagined rollouts, is accurate but pays a search cost at every step and is
sensitive to how the search is initialized. A reactive policy read out of the model is fast but
cannot look ahead. These are the fast and slow systems of control. We ask whether one world model
can hold both: a fast \emph{intuition} that proposes, and a slow deliberation that checks and
refines. We further ask whether that intuition can be pretrained once across many tasks and reused,
rather than relearned per task. GIP is our attempt at both.

\section{Method}
\textbf{Latent world model.} We build on a joint-embedding predictive architecture. An image encoder
produces a per-frame latent; an autoregressive predictor advances that latent conditioned on an
action embedding; training minimizes latent consistency (predict the next encoded frame) with a
sketched-Gaussian variance regularizer that prevents collapse. This is the LeWM base on which
everything else is mounted.

\textbf{The intuition head.} Alongside the predictor we attach a second autoregressive module, the
intuition head, that predicts the next action $a_t$ from the latent state history $z_{\le t}$,
optionally conditioned on the shifted past actions $a_{<t}$. It is trained against both the action
embedding (a JEPA-style latent target) and a decoded raw action. The head is a fast amortized
policy: one forward pass, no search.

\textbf{Intuition-guided planning.} At decision time the head rolls out a short proposed action
sequence. That proposal warm-starts CEM in place of a zero or Gaussian initialization; the world
model then rolls the candidate sequences to the goal latent and selects the best. One switch yields
three modes: \emph{reactive} (the head alone), \emph{planning} (search alone, the plain world
model), and \emph{guided} (the intuition warm-starts search).

\textbf{Generalized pretraining: padded multi-task JEPA.} To pretrain the intuition across tasks with
different action and proprioception dimensions, we pad every task to a common block and mask the
invalid dimensions before the action encoder, in the loss, and at decode. Task identity enters as a
frozen language (CLIP) token, proprioception as a separate predicted token. Task sampling uses a
temperature $\alpha$ ($p \propto n^{\alpha}$) to balance per-task effective epochs under extreme
dataset-size skew.

\section{Preliminary results}

\subsection{Intuition-guided planning matches or beats both baselines}
On goal-reaching control the reactive head alone is weak, plain planning is strong, and guiding the
planner with the intuition matches or exceeds plain planning while cutting its variance. Table~\ref{tab:modes}
reports success rate at convergence with the full CEM budget; the guided and planning columns are
paired over three seeds.

\begin{table}[h]
\centering
\small
\begin{tabular}{lccc}
\toprule
task & reactive & planning & guided \\
\midrule
PushT   & 58   & $64.7\pm6.1$ & $\mathbf{76.7\pm4.2}$ \\
TwoRoom & 44   & $\mathbf{94.7}$ & $91.3$ \\
Reacher & 2    & $\mathbf{83.3}$ & $82.7$ \\
\bottomrule
\end{tabular}
\caption{Success rate (\%) on goal-reaching. Guided wins PushT on all three seeds and ties elsewhere.
The reactive head alone is myopic on goal-reaching (Reacher 2, TwoRoom 44); its value is as a
\emph{prior for search}, not as a standalone controller.}
\label{tab:modes}
\end{table}

The benefit of guiding is robustness, not a higher ceiling: cold planning is high-variance across
seeds, whereas the guided initialization is tight, and guided reaches the same ceiling in roughly
4--5$\times$ fewer CEM iterations. On an aligned pixel manipulation task (OGBench Cube) the ordering
is sharper: guided $88 >$ planning $78 >$ the un-finetuned world-model reproduction $68$. GIP joint
training lifts planning from $68$ to $78$, and the intuition warm-start lifts $78$ to $88$.

\subsection{What the intuition encodes: contact-state memory}
Conditioning the intuition on action history helps decisively on exactly one family, contact
manipulation, and the reason is informative. On robomimic Lift, history-aware intuition raises
success from $66.0$ to $96.7$ (three seeds). A stage decomposition over 900 instrumented episodes
shows the entire gain is post-grasp: $P(\text{success}\mid\text{grasped})$ rises $76.2 \to 99.3$
while grasp attainment is unchanged. The history supplies the contact state, ``I am holding it,''
that low-resolution pixels hide. The intuition is therefore reading a latent variable the
reconstruction objective does not capture. This dictates a per-family evaluation: guided planning
for goal-reaching, history-aware reactive control for contact manipulation.

\subsection{Generalized pretraining across seven tasks}
Joint training on seven tasks with $\alpha\!\approx\!0.4$ and a doubled budget reaches or beats the
single-task reference where per-task effective epochs suffice (TwoRoom $100$ vs.\ single $93.3$;
Lift $94$; Can $66$). Uniform-frame sampling ($\alpha\!=\!1$) starves the small tasks (Square $9.3$),
confirming that the controlling quantity under size skew is effective epochs, not raw draws.
Multi-task training matches scratch single-task performance to the decimal, so the joint objective
does not interfere. Budget-matched from-scratch training also matches the warm-started endpoint, yet
a frozen-representation probe shows a large pretraining gap ($+28.7/+11.3/+16.0$ on Lift/Can/Square).
The pretrained latent is real; its benefit is efficiency and reuse rather than a higher ceiling.

\subsection{A clean, reproducible implementation}
The full method (the model, the three eval modes, and multi-task pretraining) has been re-expressed
as a standalone extension of the stable-worldmodel framework, with the intuition exposed natively
through the planner's action-prior interface. A from-scratch PushT run on that implementation
reproduces the original decode quality ($R^2\!\approx\!0.885$ vs.\ $0.88$) and the mode ordering
(guided $95 >$ planning $85 >$ reactive $65$), so the result is implementation-independent.

\section{Proposed work}
\begin{itemize}
\itemsep0.2em
\item \textbf{Scale the intuition prior.} Extend generalized intuition pretraining to more tasks and
larger data, and toward a unified action space across embodiments, padding and masking as here.
\item \textbf{Close the fast/slow loop.} Feed the planner's corrected trajectories back to retrain the
intuition, so deliberation teaches intuition and intuition shortens deliberation.
\item \textbf{Characterize when intuition helps.} Make precise where a fast prior aids search
(variance reduction, fewer iterations) versus where a noisy prior misleads it (random-action data,
e.g.\ Reacher, where guided $\le$ planning).
\item \textbf{Transfer.} Evaluate the pretrained intuition as a control prior on held-out tasks,
measuring sample- and search-efficiency against per-task learning.
\end{itemize}

\begin{thebibliography}{9}
\small
\bibitem{lewm} Q.\ Le~Lidec et al.\ \emph{Learning a world model and planning with it (LeWM).}
\bibitem{jepa} A.\ Bardes et al.\ \emph{Joint-embedding predictive architectures}; M.\ Zhou et al.\ \emph{DINO-WM}.
\bibitem{fastslow} D.\ Kahneman. \emph{Thinking, Fast and Slow}; D.\ Silver et al.\ \emph{MuZero}; policy-prior-guided search.
\bibitem{clip} A.\ Radford et al.\ \emph{Learning transferable visual models (CLIP).}
\bibitem{swm} \emph{stable-worldmodel} / \emph{stable-pretraining} frameworks.
\end{thebibliography}

\end{document}
